\documentclass[UTF8,11pt]{ctexart}
\usepackage{amsmath,amssymb,geometry}
\geometry{margin=22mm}
\begin{document}
\section*{P67: AJRP近似联合派遣与路径模型的完整路径形式F1}
\paragraph{来源与范围}依据原论文PDF第1, 7, 8, 9, 10, 11, 12, 14, 15, 16, 18, 21, 33, 34, 38, 39, 40页，第5.1节F1式(14)--(22)；第4节式(11)--(13)；第3节与附录C/F的路线定义及式(43)--(44)。

\paragraph{核心形式与范围}
保留论文第5.1节的 F1，即作者提出的近似联合调度与路径策略 AJRP 的完整路径形式。
它与第4节式(11)--(13)的当期决策问题等价，但以预估未来成本近似原随机动态 JDR；
目标值对应该近似调度模型。
下面固定当前时点 $t_n$，同时优化当期路线与其余时段的预估司机调度。

\paragraph{集合与已知数据}
时段为 $1,\ldots,N$，决策时点 $t_1<\cdots<t_N$ 给定，
当前已实现订单位置集合为 $\mathcal I^n$，配送门店为节点 $0$。
每个 $i\in\mathcal I^n$ 有需求 $q_i^n\in\mathbb Z_{>0}$，$q_i^n\le Q$；
$Q>0$ 为每条路线车辆容量。司机总数为整数 $\bar K$，车辆与司机同质。
$\zeta_p^n\in\{0,\ldots,\bar K\}$ 是在 $t_n$ 已知的、因以前时段派遣而在
$t_p$ 仍占用的司机数，$p=n,\ldots,N$。

行驶时间 $d_{ij}^n\ge0$ 对所有不同的 $i,j\in\{0\}\cup\mathcal I^n$ 给定，
它对应原文 $t_{ij}^n$，这里换记号以区别决策时点 $t_n$。
每个客户服务时间为 $s\ge0$，门店服务时间取0；订单均在 $t_n$ 已备妥
（原文取准备时间 $t_p=0$），从备妥至送达的上限为 $L_{\max}$。

对未来 $p=n+1,\ldots,N$ 与司机数量 $k=0,\ldots,\bar K$，给定离线预估的
单期配送成本 $h_p(k)$。其含义是原文
$\mathbb E_{\mathbb Q^p}[\mathcal H^s(k,\mathcal I^p,\mathbf q^p)]$，
其中 $\mathbb Q^p$ 为该期订单位置与需求的分布；本核心模型中 $h_p(k)$ 是输入系数，
若某司机数量被判为不可行而成本为 $+\infty$，
令 $\mathcal K_p=\{k:h_p(k)<+\infty\}$，固定其余 $x_k^p=0$，避免把无穷系数乘零。

给定 $\omega_m^p(k)\ge0$，表示未来 $m$ 期派遣 $k$ 名司机后，到 $p\ge m$
期仍占用的预估司机数；$\omega_m^m(k)=k$、$\omega_m^p(0)=0$。
它们按原文由历史数据或离线模拟估计，用于估计后续时段服务能力。

\paragraph{完整可行路线集合及系数}
$\mathcal R^n$ 为所有非空初等门店路线 $r=(0,i_1,\ldots,i_h,0)$ 的集合，
其中 $i_1,\ldots,i_h\in\mathcal I^n$ 两两不同。每条路线从 $t_n$ 出发、
完成整条路线后返回门店，不中途返店重派。
定义第 $\ell$ 个客户的送达时间（相对 $t_n$）和回店时长
\begin{align}
a_{r,i_1}&=d_{0i_1}^n,\nonumber\\
a_{r,i_\ell}&=a_{r,i_{\ell-1}}+s+d_{i_{\ell-1}i_\ell}^n,
&&\ell=2,\ldots,h,\tag{D1}\\
\ell_r&=a_{r,i_h}+s+d_{i_h0}^n.\tag{D2}
\end{align}
仅将满足以下条件的路线放入 $\mathcal R^n$：
\begin{equation}
\sum_{\ell=1}^h q_{i_\ell}^n\le Q,\qquad
 a_{r,i_\ell}\le L_{\max}\quad(\ell=1,\ldots,h).\tag{D3}
\end{equation}
不另外要求回店时间也小于 $L_{\max}$；它约束的是订单送达。
给每条可行路线定义固定系数
\begin{equation}
c_r=\sum_{\ell=1}^h a_{r,i_\ell},\qquad
\alpha_{i,r}=\mathbf1\{i\in r\},\qquad
\beta_{p,r}=\mathbf1\{\ell_r>t_p-t_n\},\quad p>n.\tag{D4}
\end{equation}
按原文，恰好在 $t_p$ 回店的司机可用于该期派遣，故为严格大于号。
(D1)--(D4)展开论文第8--9页、附录C及附录F式(43)--(44)对路线和送达成本的定义。
这些计算给出路线列的系数。

\paragraph{决策变量}
$\theta_r\in\{0,1\}$ 表示当期是否派一名可用司机执行路线 $r$。
$x_k^p\in\{0,1\}$ 表示未来第 $p$ 期计划派 $k$ 名司机。
$z_k^p\in\{0,1\}$ 给当期路线在未来第 $p$ 期继续占用的司机数选择一个
整数容量额度 $k$。原式(18)为上界联动，因此 $k$ 可等于或大于实际占用数；
该条件采用上述不等式。

\paragraph{唯一核心模型 F1}
\begin{align}
\min\quad&\sum_{r\in\mathcal R^n}c_r\theta_r+
\sum_{p=n+1}^N\sum_{k\in\mathcal K_p}h_p(k)x_k^p\tag{14}\\
\text{s.t.}\quad&
\sum_{m=n+1}^{p}\sum_{k=0}^{\bar K}\omega_m^p(k)x_k^m
\le\bar K-\zeta_p^n-\sum_{k=0}^{\bar K}kz_k^p,
&&p=n+1,\ldots,N,\tag{15}\\
&\sum_{k=0}^{\bar K}x_k^p=1,
&&p=n+1,\ldots,N,\tag{16}\\
&\sum_{k=0}^{\bar K}z_k^p=1,
&&p=n+1,\ldots,N,\tag{17}\\
&\sum_{r\in\mathcal R^n}\beta_{p,r}\theta_r
\le\sum_{k=0}^{\bar K}kz_k^p,
&&p=n+1,\ldots,N,\tag{18}\\
&\sum_{r\in\mathcal R^n}\theta_r\le\bar K-\zeta_n^n,\tag{19}\\
&\sum_{r\in\mathcal R^n}\alpha_{i,r}\theta_r=1,
&&i\in\mathcal I^n,\tag{20}\\
&x_k^p,z_k^p\in\{0,1\},
&&p=n+1,\ldots,N,\ k=0,\ldots,\bar K,\tag{21}\\
&\theta_r\in\{0,1\},&&r\in\mathcal R^n,\tag{22}\\
&x_k^p=0,&&p>n,\ k\notin\mathcal K_p.\tag{finite}
\end{align}
式(15)把以后派遣带来的预计占用、以前派遣带来的已知占用和本次路线的真实
跨期占用同时放入同一司机容量上限。
式(20)保证所有当期订单完整覆盖且仅一次，没有拆单或推迟到下一波。

\paragraph{原文消歧与边界}
原 PDF 式(15)的 $\zeta$ 上下标排成 $\zeta_n^{p}$；按第8页状态定义、
式(12)及附录A，统一为 $\zeta_p^n$，即当前已知的未来占用数。
$n=N$ 时未来求和为空，模型自然保留当期完整配送路线选择。
$\mathcal I^n=\varnothing$ 时没有非空当期路线，司机可以不派出；
未来司机数量0只在输入成本允许时使用，不无依据强制每期至少派一人。
本模型采用全部可行初等路线；仅使用任意一小部分候选路线会得到受限模型，
其最优值对应所列候选路线构成的受限模型。

\end{document}
